\documentclass{article}

% ============================================================
% MA 213 In-Class Worksheet Template
% Student-facing worksheet for lecture activities.
% ============================================================

\usepackage[margin=1in]{geometry}
\usepackage{amsmath}
\usepackage{xcolor}
\usepackage{parskip}
\usepackage{array}
\usepackage{enumitem}
\usepackage{booktabs}

\newcommand{\coursenumber}{MA 213}
\newcommand{\weekplaceholder}{4}
\newcommand{\lectureplaceholder}{8}
\newcommand{\lecturetitle}{Bayes' Theorem: A Medical Test: Think/Pair/Share}

\newcommand{\nameblank}{\rule{1.75in}{0.4pt}}
\newcommand{\idblank}{\rule{1.55in}{0.4pt}}
\newcommand{\shortblank}{\rule{1in}{0.4pt}}
\newcommand{\longblank}{\rule{0.93\linewidth}{0.4pt}}

\begin{document}

\begin{center}
  {\LARGE \coursenumber{} Week \weekplaceholder{}: Lecture \lectureplaceholder{}}\\
  {\large \lecturetitle{}}
\end{center}

\begin{enumerate}[label=\arabic*.,leftmargin=*,itemsep=.7em]
  \item \textbf{Your name:} \nameblank \hfill \textbf{BU ID:} \idblank
\end{enumerate}

\textbf{Big idea:} We just used the Law of Total Probability and Bayes' Theorem to work out a weather-prediction example. Now apply the exact same method to a new, more consequential context: interpreting a medical test.

\section*{The Setup}

A medical test is given for a rare disease.
\begin{itemize}
  \item The disease occurs in 1\% of the population.
  \item If a person has the disease, the test is positive 95\% of the time.
  \item If a person does not have the disease, the test is positive 5\% of the time.
\end{itemize}
Suppose a randomly chosen person tests positive.

\section*{Step 1: Think (2 mins)}
Hint: this has the exact same structure as the weather-prediction example --- there are more hints on the next page if you get stuck.

\textbf{(a)} Define events $A_1 = \{\text{has the disease}\}$, $A_2 = \{\text{does not have the disease}\}$, and $B = \{\text{tests positive}\}$. Write down $P(A_1)$, $P(A_2)$, $P(B \mid A_1)$, and $P(B \mid A_2)$.
\vspace{0.6in}

\textbf{(b)} Use the Law of Total Probability to compute $P(B) = P(\text{Positive})$.
\vspace{0.6in}

\textbf{(c)} Use Bayes' Theorem to compute $P(A_1 \mid B) = P(\text{Disease} \mid \text{Positive})$.
\vspace{0.6in}

\newpage
\section*{Hints for Step 1}

\textbf{(a)} Match each given percentage to one of $P(A_1)$, $P(A_2)$, $P(B\mid A_1)$, $P(B\mid A_2)$. Remember $A_2$ is the complement of $A_1$, so $P(A_2) = 1 - P(A_1)$.

\vspace{0.2in}

\textbf{(b)} Recall the formula from lecture:
\[ P(B) = P(B\mid A_1)P(A_1) + P(B\mid A_2)P(A_2) \]
Plug in the four numbers from part (a).

\vspace{0.2in}

\textbf{(c)} Recall Bayes' Theorem:
\[ P(A_1 \mid B) = \frac{P(A_1)P(B\mid A_1)}{P(B)} \]
You already computed the denominator in part (b) --- reuse it.

\vspace{0.2in}
\noindent\rule{\linewidth}{0.4pt}
\vspace{0.2in}

\section*{Step 2: Pair (2 mins)}

\begin{enumerate}[label=\arabic*.,leftmargin=*,itemsep=.7em, start=2]
  \item \textbf{Partner's name:} \nameblank \hfill \textbf{BU ID:} \idblank
\end{enumerate}

Compare your answers with your partner. Then discuss:

\begin{itemize}
  \item Is $P(\text{Disease} \mid \text{Positive})$ higher or lower than you expected, given that the test is ``95\% accurate''?
  \item The test has the same 95\%/5\% accuracy either way, but the disease is rare. How does the rarity of the disease affect the answer?
  \item What would happen to $P(\text{Disease} \mid \text{Positive})$ if the disease were much more common, say 50\% of the population? (You don't need to recompute --- just reason about it.)
\end{itemize}

\textbf{Our thoughts:}
\vfill

\end{document}
